\documentclass[10pt,a4paper]{amsart}
\usepackage[margin=19mm]{geometry}
\usepackage{amsmath,amssymb,mathtools}
\usepackage[scr=rsfs,cal=euler]{mathalfa}
\usepackage{xfrac}
\usepackage[unicode,hidelinks,bookmarks=false]{hyperref}
\newcommand{\ie}{i.e.\ }
\newcommand{\cf}{cf.\ }
\newcommand{\resp}{respectively\ }
\newcommand{\Subsection}{\subsection}
\providecommand{\nobreakdash}{\nobreak\hbox{-}}
\setlength{\emergencystretch}{2em}
\multlinegap0pt
\usepackage{bm}
\usepackage{bbm}
\usepackage{dsfont}
\usepackage{xspace}
\usepackage{cancel}
\usepackage{mathtools}
\usepackage{amsthm}
\makeatletter
\@ifundefined{theorem}{\newtheorem{theorem}{Theorem}}{}
\@ifundefined{lemma}{\newtheorem{lemma}[theorem]{Lemma}}{}
\makeatother
\makeatletter
\newcommand{\FF}{\mathbb{F}}
\expandafter\ifx\csname caseof\endcsname\relax
\newenvironment{caseof}{\setcounter{casenum}{0}}{\vskip.5\baselineskip}
\fi
\expandafter\ifx\csname caseofb\endcsname\relax
\newenvironment{caseofb}{\setcounter{casenumb}{0}}{\vskip.5\baselineskip}
\fi
\renewcommand{\le}{\leqslant}
\renewcommand{\leq}{\leqslant}
\newcommand{\one}{\boldsymbol{1}}
\expandafter\ifx\csname subcaseof\endcsname\relax
\newenvironment{subcaseof}{\setcounter{subcasenum}{0}}{\vskip.5\baselineskip}
\fi
\expandafter\ifx\csname subcaseofb\endcsname\relax
\newenvironment{subcaseofb}{\setcounter{subcasenumb}{0}}{\vskip.5\baselineskip}
\fi
\makeatother
\title[Quantum Codes with Addressable and Transversal Non-Clifford ]{Quantum Codes with Addressable and Transversal Non-Clifford Gates}
\author{Zhiyang He, Vinod Vaikuntanathan, Adam Wills, Rachel Yun Zhang}
\date{Source: arXiv:2502.01864v3 (CC BY 4.0)}
\begin{document}
\maketitle

\noindent\emph{Source and licence.} Quantum Codes with Addressable and Transversal Non-Clifford Gates. Version arXiv:2502.01864v3 (CC BY 4.0), distributed under the Creative Commons Attribution 4.0 International licence, as recorded in the arXiv licence metadata for this exact version and shown on the abstract page https://arxiv.org/abs/2502.01864v3. Full rights notice: licenses/arxiv-2502.01864-CC-BY-4.0.txt.\par\smallskip

\noindent\textit{Editorial scope.} Counted unit: the Lemma whose source label is \texttt{lem:extended\_orthogonality} in the article (source0.tex lines 1392--1411, source label \texttt{lem:extended\_orthogonality}), delivered together with the article's own complete original proof of it (source0.tex lines 1412--1502, one \texttt{proof} environment of 91 lines). No mathematics is written, repaired or re-derived: statement and proof are the article's own text line for line. Required context retained uncounted: eq:matrix\_M (lines 1360--1367). Every definition, display and label of those lines is retained unchanged. Omitted from this package and not redistributed: 1--1359, 1368--1391, 1503--2315. Nothing inside a retained excerpt is omitted or altered: every retained body is delivered exactly as the article's lines are written, so no omission receipt of any kind accompanies this package. Citations: the retained excerpts carry no citation command at all, so no bibliography entry of the article is transcribed and no reference is redistributed. Reference adaptation: none was needed and none was made; every reference inside the delivered unit resolves inside it, so no unresolved reference or citation is produced. Printed slips kept literal: no printed word, symbol, hypothesis or number of the counted statement or of its proof was altered. Layout and class: the article's own class is \texttt{article} with packages including hyperref, geometry, mathptmx, amssymb, amsmath, amsthm, amsfonts, fontenc, bm, color, comment, cleveref, xcolor, float; the wrapper is the lane's pinned \texttt{amsart} editorial wrapper at 10pt on A4, which restates the theorem-like environments the retained text needs and adds the article's own macro definitions that the retained text needs (\texttt{\textbackslash FF}, \texttt{\textbackslash le}, \texttt{\textbackslash leq}, \texttt{\textbackslash one}), with extra wrapper packages (bm, bbm, dsfont, xspace, cancel, mathtools). The delivered numbering is the unit's own. Assets: no retained span contains a figure environment, a graphics-inclusion command, a PDF-page inclusion, a table or a TikZ picture; no image file is copied into this package, the package declares no asset dependency and no image file is redistributed with it. Licence: version arXiv:2502.01864v3 of this article is distributed under the Creative Commons Attribution 4.0 International licence, as recorded in the arXiv licence metadata for this exact version and shown on the abstract page https://arxiv.org/abs/2502.01864v3. A full rights notice is delivered at licenses/arxiv-2502.01864-CC-BY-4.0.txt. Characters: every retained excerpt is pure ASCII, so the delivered engine, XeLaTeX with its default Latin Modern text font, reports no missing character.
\begin{equation}\label{eq:matrix_M}
    M = 
    \begin{bmatrix}
        1 & 1 & 1 & 0 & 1 & 0 & 0 \\
        1 & 1 & 0 & 1 & 0 & 1 & 0 \\
        1 & 0 & 1 & 1 & 0 & 0 & 1
    \end{bmatrix}.
\end{equation}

\begin{lemma}\label{lem:extended_orthogonality}
    For $0\le A\le k'-1$, let $\Lambda^A = \sum_{j=1}^7 \Gamma^{7A+j}$. 
    Let $T(x_1, x_2, x_3) = x_1^{D_1}x_2^{D_2}x_3^{D_3}$ be a degree-7 monomial on three variables, i.e., $D_1 + D_2 + D_3 = 7$.
    For all $a, b, c\in [m']$, if $D_1,D_2,D_3$ are all non-zero, 
    we have
    \begin{align}
        \label{eq:intrablock_lem_1}
        \sum_{i=1}^n \Lambda^A_iT(h^{a}_i, h^b_i,h^c_i) = 
        \begin{cases}
            1 &\text{ if } \{a, b, c\} = \{3A+1, 3A+2, 3A + 3\}, \\
            0 &\text{ otherwise.}
        \end{cases}
    \end{align}
    If any of $D_1,D_2,D_3$ are zero, 
    we have
    \begin{align}
    \label{eq:intrablock_lem_2}
        \sum_{i=1}^n \Lambda^A_iT(h^{a}_i, h^b_i,h^c_i) = 0.
    \end{align}
\end{lemma}

\begin{proof}
    Let us expand the summation as 
    \begin{align}
        \sum_{i=1}^n \Lambda^A_iT(h^{a}_i, h^b_i,h^c_i)
        &= 
        \sum_{i=1}^n \sum_{j=1}^7 \Gamma^{7A+j}_i T(h^{a}_i, h^b_i,h^c_i) \\
        &= 
        \sum_{j=1}^7 \left(\sum_{i=1}^n \Gamma^{7A+j}_i T(h^{a}_i, h^b_i,h^c_i)\right) \\
        &=\sum_{j=1}^7 \left(\sum_{i=1}^n \Gamma^{7A+j}_i (h^{a}_i)^{D_1} (h^b_i)^{D_2}(h^c_i)^{D_3}\right).\label{eq:calculation_designed_intrablock_lemma}
    \end{align}
    The rows of $H$ are linear combinations of rows of $G$. 
    Recall the matrix $M$ in equation~\eqref{eq:matrix_M},
    we denote its rows as $M^1, M^2, M^3$, so that its $(i,j)$-th entry is $M^i_j$. 
    For $0\le B_a\le k'-1$, for $e_a\in [3]$,
    \begin{align}
        h^{3B_a+e_a} &= \sum_{j_a=1}^7 M^{e_a}_{j_a} g^{7B_a+j_a}.
    \end{align}
    For all $a > 3k'$, we simply have $h^a = g^{a + 4k'}$, as mentioned.
    
    We will now do a case analysis. Note that Equation~\eqref{eq:calculation_designed_intrablock_lemma} will be evaluated using the fact that $h^a$ are either linear combinations of rows of $G_1$ (if $a \leq 3k'$), or are simply rows of $G_0$ (if $a > 3k'$). It is therefore easy to see that if any of $a,b,c$ are greater than $3k'$, this sum must vanish using the addressable 7-orthogonality of $G$. For example suppose that $a, b, c > 3k'$ and all $D_1, D_2, D_3$ are non-zero. Then,
    \begin{align}
        \sum_{j=1}^7 \left(\sum_{i=1}^n \Gamma^{7A+j}_i (h^{a}_i)^{D_1} (h^b_i)^{D_2}(h^c_i)^{D_3}\right)
        &= \sum_{j=1}^7 \left(\sum_{i=1}^n \Gamma^{7A+j}_i 
        \left(g^{a+4k'}_i \right)^{D_1} 
        \left(g^{b+4k'}_i \right)^{D_2}
        \left(g^{c+4k'}_i \right)^{D_3}
        \right) \\
        &= \sum_{j=1}^7 \one_{a+4k'=b+4k'=c+4k'=7A+j} \label{eq:intrablock_lem_calc_indicator} \\
        &= \sum_{j=1}^7 0 = 0.
    \end{align}
    Note that equation~\eqref{eq:intrablock_lem_calc_indicator} is an application of the addressable $7$-orthogonality property of $G$.
    Moreover, if $D_3 = 0$, we simply replace the indicator function in equation~\eqref{eq:intrablock_lem_calc_indicator} with $\one_{a+4k'=b+4k'=7A+j}$ and the same result hold. Same is true if any of $D_1,D_2,D_3$ is zero. 
     
    The other cases, where we only assume that at least one of $a,b,c$ is more than $3k'$, follow for the same reason.

    The only non-trivial case is therefore $a,b,c \leq 3k'$. To this end, let us show Equations~\eqref{eq:intrablock_lem_1} and~\eqref{eq:intrablock_lem_2} with $a = 3B_a + e_a, b = 3B_b + e_b$, and $c = 3B_c + e_c$. 
    Beginning with the case that $D_1, D_2$ and $D_3$ are all non-zero, Equation~\eqref{eq:calculation_designed_intrablock_lemma} evaluates as follows. 
    \begin{align}
        &\sum_{j=1}^7 \left(\sum_{i=1}^n \Gamma^{7A+j}_i (h^{a}_i)^{D_1} (h^b_i)^{D_2}(h^c_i)^{D_3}\right) \\
        &= \sum_{j=1}^7 \left(\sum_{i=1}^n \Gamma^{7A+j}_i 
        \left(\sum_{j=1}^7 M^{e_a}_{j} g^{7B_a+j}_i \right)^{D_1} 
        \left(\sum_{j=1}^7 M^{e_b}_{j} g^{7B_b+j}_i \right)^{D_2} 
        \left(\sum_{j=1}^7 M^{e_c}_{j} g^{7B_c+j}_i \right)^{D_3}
        \right) \\
        &= \sum_{j=1}^7 
        \sum_{\substack{J_a\in [7]^{D_1}}}
        \sum_{\substack{J_b\in [7]^{D_2}}}
        \sum_{\substack{J_c\in [7]^{D_3}}}
        M^{e_a}_{J_a}M^{e_b}_{J_b}M^{e_c}_{J_c}
        \left(
        \sum_{i=1}^n \Gamma^{7A+j}_i g_i^{7B_a + J_a}g_i^{7B_b + J_b} g_i^{7B_c + J_c} \right) \label{eq:intrablock_lem_bigsum}
    \end{align}
    Here, we have introduced a short-hand to more easily expand products like $\left(\sum_{j=1}^7M_j^{e_a}g_i^{7B_a+j}\right)^{D_1}$. Indeed, given a tuple $J\in [7]^D$, we denote $M^e_J = \prod_{\delta = 1}^D M^e_{J_\delta}$ and $g^{7B+J}_i = \prod_{\delta = 1}^D g^{7B + J_\delta}_i$.
    Again by the addressable $7$-orthogonality property of $G$, $\sum_{i=1}^n \Gamma^{7A+j}_i g_i^{7B_a + J_a}g_i^{7B_b + J_b} g_i^{7B_c + J_c}$ is 1 when $B_a = B_b = B_c = A$, and when all the indices of $J_a, J_b$ and $J_c$ are equal to $j$.
    Therefore, equation~\eqref{eq:intrablock_lem_bigsum} simplifies to
    \begin{align}
        \sum_{j=1}^7 
        (M^{e_a}_{j})^{D_1} (M^{e_b}_{j})^{D_2} (M^{e_c}_{j})^{D_3}.
        \label{eq:intrablock_lem_mult_of_M}
    \end{align}
    In fact, one can check that Equation~\eqref{eq:calculation_designed_intrablock_lemma} evaluates to~\eqref{eq:intrablock_lem_mult_of_M} whenever $a=3B_a+e_a$, $b=3B_b+e_b$ and $c=3B_c+e_c$, irrespective of whether any of $D_1, D_2, D_3$ are zero.
    Let us now inspect the matrix $M$. 
    Observe that the entries in $M$ are all $0$ and $1$, which means 
    \begin{align}
        (M^{e_a}_j)^{D_1} = 
        \begin{cases}
            M^{e_a}_j &\text{ if $D_1\ne 0$,}\\
            1         &\text{ if } D_1 = 0.
        \end{cases}
        \label{eq:intrablock_lem_binary_exp}
    \end{align}
    Moreover, the three rows of $M$ satisfy the following properties: 
    \begin{align}
        \sum_{j=1}^7 M^{e}_j &= 0 \text{ for all $e\in [3]$, } \\
        \sum_{j=1}^7 (M^{e_1}\star M^{e_2})_j &= 0 \text{   for all $e_1,e_2\in [3]$,} \\
        \sum_{j=1}^7 (M^{e_1}\star M^{e_2}\star M^{e_3})_j &=
        \begin{cases}
            1 &\text{ if $\{e_1,e_2,e_3\} = \{1,2,3\}$,} \\
            0 &\text{ otherwise.}
        \end{cases}
        \label{eq:intrablock_lem_mult_prop_M}
    \end{align}
    We note that all summations in the above three equations are over the field $\FF_q$, and the above equations hold because $\FF_q$ has characteristic $2$.
    Combining equations~\eqref{eq:intrablock_lem_binary_exp} to~\eqref{eq:intrablock_lem_mult_prop_M}, we see that 
    \begin{align}
        \sum_{j=1}^7 
        (M^{e_a}_{j})^{D_1} (M^{e_b}_{j})^{D_2} (M^{e_c}_{j})^{D_3} = 1
    \end{align}
    if and only if $D_1,D_2,D_3$ are non-zero, and $\{e_a,e_b,e_c\} = \{1,2,3\}$. 
    This implies that $\{a,b,c\} = \{3A+1,3A+2,3A+3\}$, as desired. 
\end{proof}

\end{document}
